\section{Modélisation Deep Q-Network (DQN)}
\label{sec:dqn}

Pour faire suite aux développements initiaux basés sur l'algorithme \textit{Proximal Policy Optimization} (PPO) détaillés précédemment, l'approche a été étendue à l'algorithme \textbf{Deep Q-Network (DQN)} afin d'étudier l'impact d'une méthode de type valeur (\textit{value-based}) vis-à-vis d'une méthode de gradient de politique (\textit{policy-gradient}) sur le compromis énergie/latence au sein de notre infrastructure \textit{Smart City}.

Contrairement à PPO, qui modélise nativement des distributions de probabilités sur des espaces d'actions multidimensionnels indépendants via \texttt{MultiDiscrete}, DQN évalue la qualité de chaque action possible au travers d'un réseau de neurones estimateur de la fonction de valeur de sous-action optimale $Q(s, a)$. Cette spécificité algorithmique a imposé une refonte complète de la représentation de l'espace d'action et du cycle de traitement des décisions.

\subsection{Standardisation et Discrétisation de l'Environnement}
L'interface standardisée \texttt{gymnasium} sert de passerelle avec le simulateur à événements discrets \textit{FaaSim/Skippy}. La discrétisation temporelle conserve une logique de pause périodique toutes les $X$ secondes (\texttt{step\_duration}).

L'espace d'observation reste inchangé afin de garantir une comparaison équitable entre les modèles. L'agent observe un espace continu matérialisé par une \texttt{spaces.Box} qui encapsule :
\begin{itemize}
    \item La latence moyenne d'exécution des dernières requêtes par fonction.
    \item Le nombre exact d'instances (\texttt{replicas}) actives pour chacune des 6 fonctions déployées.
\end{itemize}

\subsection{Linéarisation et Décodage de l'Espace d'Action}
La contrainte majeure de DQN réside dans son incapacité native à gérer des espaces d'actions combinatoires ou multidimensionnels de type \texttt{MultiDiscrete}. Pour y remédier, un protocole de linéarisation (\textit{flattening}) a été développé au sein de la classe \texttt{SmartCityGymEnvDQN}.

Le système comporte 6 fonctions applicatives indépendantes, possédant chacune 17 ordres de grandeur de \textit{scaling} possibles (de $-8$ pour une réduction maximale à $+8$ pour une augmentation maximale, l'action $0$ désignant le maintien de l'état). L'espace d'action DQN est donc configuré comme un espace discret unique :
\begin{equation}
    a_{\text{DQN}} \in \mathcal{A} = \text{\texttt{spaces.Discrete}}(6 \times 17) = \text{\texttt{spaces.Discrete}}(102)
\end{equation}

Lors de chaque interaction (\texttt{step}), l'action scalaire émise par le réseau de neurones DQN est projetée mathématiquement pour identifier la fonction cible et la valeur brute de scaling associée :
\begin{equation}
    \text{func\_idx} = \lfloor a_{\text{DQN}} / 17 \rfloor
\end{equation}
\begin{equation}
    \text{act\_raw} = a_{\text{DQN}} \pmod{17}
\end{equation}
\begin{equation}
    \text{act\_val} = \text{act\_raw} - 8
\end{equation}

\begin{figure}[htbp]
\begin{center}
\begin{tikzpicture}[node distance=2cm, auto, >=latex']
    \node [draw, rectangle, fill=blue!10, rounded corners, minimum width=3cm, minimum height=1cm] (input) {Action DQN unique ($a \in [0, 101]$)};
    \node [draw, rectangle, fill=gray!10, minimum width=2.5cm, minimum height=1cm, below left of=input, node distance=3.5cm] (div) {Division entière : $\lfloor a / 17 \rfloor$};
    \node [draw, rectangle, fill=gray!10, minimum width=2.5cm, minimum height=1cm, below right of=input, node distance=3.5cm] (mod) {Modulo : $a \pmod{17}$};
    \node [draw, circle, fill=green!10, below of=div, node distance=2.5cm] (func) {Fonction Cible (0 à 5)};
    \node [draw, rectangle, fill=green!10, below of=mod, node distance=2.5cm] (val) {Ordre Réel : $\text{act\_raw} - 8$};
    
    \draw[->, thick] (input) -- (div.north);
    \draw[->, thick] (input) -- (mod.north);
    \draw[->, thick] (div) -- (func);
    \draw[->, thick] (mod) -- (val);
    
    \node[draw, dashed, rectangle, fit=(func) (val), inner sep=0.3cm, label=below:Actionneur du Simulateur FaaSim] (actuator) {};
\end{tikzpicture}
\end{center}
\caption{Schéma de décodage de l'espace d'action linéarisé pour l'algorithme DQN.}
\label{fig:action_decoding}
\end{figure}

Le listing \ref{lst:dqn_step} illustre l'implémentation de ce décodage et l'injection des requêtes de modification d'état asynchrones dans le moteur de simulation \textit{SimPy}.

\begin{lstlisting}[caption=Décodage de l'action et logique d'actuation au sein du step DQN, label=lst:dqn_step]
# Decodage de l'action scalaire globale emise par le DQN
func_idx = int(action) // 17
act_raw = int(action) % 17
act_val = act_raw - 8
fname = self.functions[func_idx]
current_replicas = self.get_current_replicas(fname)

if act_val < 0:
    pods_to_remove = abs(act_val)          
    if current_replicas >= pods_to_remove: 
        self.sim_env.process(self.sim_env.faas.scale_down(fname, pods_to_remove))
    elif current_replicas > 0:
        self.sim_env.process(self.sim_env.faas.scale_down(fname, current_replicas))
elif act_val > 0: 
    if current_replicas + act_val <= 15: 
        self.sim_env.process(self.sim_env.faas.scale_up(fname, act_val))
    elif current_replicas < 15: 
        pods_to_add = 15 - current_replicas
        self.sim_env.process(self.sim_env.faas.scale_up(fname, pods_to_add))
# Si act_val == 0, l'environnement reste stable (statu quo)
\end{lstlisting}

\subsection{Hyperparamètres, Exploration et Suivi de l'Apprentissage}
L'architecture DQN repose sur un réseau de neurones profond couplé à un mécanisme de \textit{Replay Buffer} (mémoire d'expérience) pour briser les corrélations temporelles. L'apprentissage est piloté via la bibliothèque \texttt{stable-baselines3}.

Contrairement à PPO qui explore par perturbation stochastique de sa politique, DQN utilise une stratégie $\epsilon$-greedy. L'exploration décroît de manière linéaire via la configuration du paramètre \texttt{exploration\_fraction}. Le suivi de la convergence et des performances d'infrastructure est assuré par une stack \textbf{TensorBoard} dédiée stockée sous \texttt{./logs\_tensorboard\_dqn/}, permettant l'analyse des courbes de perte de la fonction de coût Bellman (\textit{loss}) et de l'évolution de la récompense par épisode.

\subsection{Modélisation Rigoureuse de la Récompense Multi-Objectif}
La fonction de récompense est calquée sur une approche normalisée strictement positive pour préserver la cohérence comparative avec le modèle PPO. Elle formule explicitement le compromis performance-efficience :
\begin{equation}
    \mathcal{R} = (\alpha \times \text{SLA}) + (\beta \times \text{Energy}) + (\gamma \times \text{Efficiency})
\end{equation}
Sous la contrainte stricte $\alpha + \beta + \gamma = 1$, fixant la borne supérieure théorique de la récompense à $1.0$.

\begin{itemize}
    \item \textbf{Score SLA} : Évalue le respect des budgets de latence. Si la latence moyenne mesurée pour une fonction franchit le seuil fixé dans le fichier de configuration global \texttt{config.json}, la note décroît linéairement vers 0.
    \item \textbf{Score Energy} : Afin d'empêcher l'agent d'exploiter les biais de calcul en simulant des maximums artificiels via l'activation de nœuds vides, la puissance de référence intègre l'inventaire matériel total physique :
    \begin{equation}
        P_{\text{max}} = f(\text{\texttt{self.sim\_env.cluster.nodes}})
    \end{equation}
    Toute activation d'infrastructure engendre ainsi une consommation réelle immédiate agissant comme une pénalité énergétique incontournable.
    \item \textbf{Score Efficiency} : Mesure le ratio $\frac{\text{Invocations}}{\text{Pods Actifs}}$ pour éliminer le sur-provisionnement logiciel (\textit{over-provisioning}) induit par la consommation statique (\textit{static power}) des serveurs Edge.
\end{itemize}

\subsection{Co-Pilotage du Placement Matériel (Hardware Policy)}
Afin de pallier les limitations induites par le placement aveugle des pods par le planificateur par défaut, l'agent DQN s'interface de manière non-intrusive avec le pipeline de priorité du \textit{Scheduler} externe via la classe \texttt{HardwarePolicyPriority}.

Bien que l'action DQN extraite du réseau de neurones pilote en premier lieu le volume de pods (Listing \ref{lst:dqn_step}), le choix de la politique matérielle sous-jacente (Modes : \textit{Éco}, \textit{Normal}, \textit{Performance}) est échantillonné de façon conjointe ou via des règles d'état heuristiques. Si le mode \textit{Éco} est sélectionné, un bonus de 1000 points est injecté dynamiquement sur les nœuds basse consommation de la topologie (ex: \textit{Raspberry RockPi}, \textit{Nvidia Jetson Nano}) pour forcer le planificateur à privilégier la sobriété énergétique lors du déploiement physique des conteneurs.